\begin{tikzpicture}[x=1cm,y=1cm,every node/.append style={anchor=center,inner sep=.3333em}]
\draw[-{Latex[length=2mm]},Ink,line width=.6pt] (-3.4,0)--(3.6,0) node[pNote,anchor=north east]{$\eta=\log\frac{p}{1-p}$};
\draw[-{Latex[length=2mm]},Ink,line width=.6pt] (0,-.2)--(0,3.1) node[pNote,anchor=south]{$p$};
\draw[Muted,dashed] (-3.3,2.6)--(3.3,2.6);
\node[pNote,anchor=east] at(-.05,2.6){1};
\draw[GeoA,line width=1.4pt,domain=-3.3:3.3,samples=80] plot(\x,{2.6/(1+exp(-1.3*\x))});
\fill[GeoB] (0,1.3) circle(2.2pt) node[pLabel,anchor=west,text=GeoB]{\ $p=0{,}5$};
\node[pNote,text width=7.3cm] at(0,-1.25){Il logit porta l’intervallo $(0,1)$ sull’intera retta. Un coefficiente $\beta$ moltiplica le probabilità a favore, $p/(1-p)$, per $e^{\beta}$ a ogni unità del predittore.};
\end{tikzpicture}
